% Glosario — términos técnicos y siglas [REGLAS_FORMATO §3]. Origen: v3.0.
% Mantener sincronizado con 00_proyecto/GLOSARIO.md. Orden alfabético.
\chapter*{Glosario}
\addcontentsline{toc}{chapter}{Glosario}

\begin{glosario}
  \item[ACV] Análisis de Ciclo de Vida.
  \item[ADT] Altura Dinámica Total. Energía requerida por un sistema de bombeo para transportar un fluido.
  \item[AGR1] Punto de muestreo correspondiente a los residuos agroquímicos generados en la bodega.
  \item[AGR2] Punto de muestreo correspondiente a los residuos agroquímicos captados en la sala de bombas.
  \item[ANID] Agencia Nacional de Investigación y Desarrollo.
  \item[AS1] Punto de muestreo correspondiente a las aguas servidas provenientes del sector de oficinas.
  \item[AS2] Punto de muestreo correspondiente a las aguas servidas habitacionales provenientes de la casa.
  \item[ASD] Diseño por Tensiones Admisibles (\textit{Allowable Stress Design}).
  \item[DBO$_5$] Demanda Bioquímica de Oxígeno medida a los cinco días.
  \item[DGA] Dirección General de Aguas.
  \item[DOH] Dirección de Obras Hidráulicas.
  \item[DQO] Demanda Química de Oxígeno.
  \item[EETT] Especificaciones Técnicas.
  \item[EST] Aguas del estero Batuco.
  \item[HC] Humedal Construido.
  \item[HC-FV] Humedal Construido de Flujo Subsuperficial Vertical.
  \item[IBC] Contenedor Intermedio para Graneles (\textit{Intermediate Bulk Container}).
  \item[INN] Instituto Nacional de Normalización.
  \item[LRFD] Diseño por Factores de Carga y Resistencia (\textit{Load and Resistance Factor Design}).
  \item[NMP] Número Más Probable (recuento de coliformes).
  \item[NT] Nitrógeno Total.
  \item[NTK] Nitrógeno Total Kjeldahl. Suma del nitrógeno orgánico y amoniacal.
  \item[OD] Demanda de Oxígeno del lecho filtrante (método del balance de oxígeno de Platzer). El oxígeno disuelto se escribe sin sigla.
  \item[OI] Aporte de Oxígeno al lecho filtrante.
  \item[PEAD] Polietileno de Alta Densidad.
  \item[PT] Fósforo Total.
  \item[PVC] Policloruro de Vinilo.
  \item[SbN] Soluciones Basadas en la Naturaleza.
  \item[SST] Sólidos Suspendidos Totales.
  \item[TRH] Tiempo de Retención Hidráulica.
\end{glosario}
